\chuongkhongso{TÓM TẮT}

Báo cáo trình bày \textbf{ScholarLens}, hệ thống truy xuất thông tin và hỏi đáp cho bài báo khoa học trên arXiv.
Người dùng tìm bài bằng câu hỏi tự nhiên (tiếng Việt hoặc tiếng Anh), thêm tài liệu bằng cách kéo-thả file PDF,
và trò chuyện với chatbot. Chatbot trả lời có trích dẫn tới đúng bài báo, kèm các bài liên quan tìm được trên đồ thị
tri thức.

Hệ thống gồm bảy module chạy độc lập: thu thập dữ liệu từ arXiv API; nạp PDF với ba lớp kiểm tra trùng (băm
SHA-256, mã arXiv, tiêu đề gần giống); embedding đa ngôn ngữ bằng \code{multilingual-e5-base} (tự dùng GPU); đồ thị
tri thức nối bài báo với tác giả, khái niệm, chuyên mục, quan hệ trích dẫn và tương đồng nội dung; chatbot hai chế
độ dựa trên Graph RAG; máy chủ LLM mã nguồn mở Qwen2.5 chạy bằng vLLM trên GPU Kaggle, công khai qua Cloudflare
Tunnel; và web app FastAPI có giao diện chat tối giản. Dữ liệu được tổ chức thành bốn kho: SQLite (metadata),
ChromaDB (vector, chỉ mục HNSW), chỉ mục BM25 và đồ thị NetworkX. Tìm kiếm kết hợp BM25 và vector bằng Reciprocal
Rank Fusion.

Trên kho thử nghiệm 492 bài, tìm bài theo tiêu đề bằng BM25 đạt Recall@1 = 0,957; cơ chế kiểm tra trùng phát hiện
95\% tiêu đề bị biến đổi mà không báo nhầm bài thật nào; trong năm bài liên quan mà đồ thị gợi ý, 37\% cùng chuyên
mục với bài gốc, gấp 2,7 lần so với chọn ngẫu nhiên.

\textbf{Từ khóa:} Information Retrieval, Hybrid Search, Graph RAG, Knowledge Graph, Chatbot, vLLM, arXiv.
